\documentclass{article}
\usepackage[utf8]{inputenc}

\usepackage{amsfonts}
\usepackage{amssymb}
\usepackage{amsmath}
\usepackage{amsthm}
\usepackage{enumitem}

\usepackage{bbold}
\usepackage{bm}
\usepackage{graphicx}
\usepackage{color}
\usepackage{hyperref}
\usepackage[margin=2.5cm]{geometry}

\begin{document}

% ==============================================================================

\title{\Large{INFO8006: Project 2 - Report}}
\vspace{1cm}
% TODO: replace with your own name(s) and student number(s)
\author\small{\bf Mr Clark Anderson Esparagoza - B6739300} 

\maketitle

% ==============================================================================

\section{Bayes filter}

\begin{enumerate}[label=\alph*.,leftmargin=*]
    \item \textbf{Sensor model.}
    Let $x_t=(w,h)$ be the position of a ghost, $q_t$ the position of Pacman
    and $d_t(x)=|w-q_{t,1}|+|h-q_{t,2}|$ the Manhattan distance between a
    cell $x$ and Pacman. The sensor returns the noisy distance
    \begin{equation*}
        E_t = d_t(x_t) + N_t, \qquad N_t = B_t - np, \qquad
        B_t \sim \mathrm{Bin}(n,p),
    \end{equation*}
    with $p=\tfrac12$ and $n=\sigma^2/(p(1-p))=4\sigma^2$ (rounded down to an
    integer), where $\sigma^2$ is the sensor variance (default $\sigma^2=1$,
    hence $n=4$). The noise $N_t$ has mean $0$ and variance
    $np(1-p)=\sigma^2$; it is independent of the ghost position and drawn
    independently at each time step. Therefore, for an observed value $e$,
    \begin{equation*}
        P(E_t=e \mid X_t=x, q_t) =
        \begin{cases}
            \binom{n}{k}\,p^{k}(1-p)^{n-k}
                & \text{if } k := e - d_t(x) + np \in \{0,\dots,n\},\\
            0 & \text{otherwise.}
        \end{cases}
    \end{equation*}
    In words, the measurement can only lie within $\pm np$ of the true
    distance, so a measurement is compatible with a thick diamond-shaped ring
    of cells around Pacman, whose thickness grows with $\sigma^2$.

    \item \textbf{Unified transition model.}
    Let $\mathcal N(x)$ be the set of the cells adjacent to $x$ that are not
    walls. A ghost never stays still and, in this project, may move backward,
    so it moves to one of the cells of $\mathcal N(x)$. Reading the three ghost
    implementations, the probability of each move is proportional to a weight:
    $1$ for every move for the \texttt{confused} ghost; $2$ if the move does
    not bring the ghost closer to Pacman and $1$ otherwise for the
    \texttt{afraid} ghost; and $2^3=8$ instead of $2$ for the \texttt{scared}
    ghost. These three ghosts are therefore obtained from a single model with
    one free parameter $\theta\ge 1$:
    \begin{equation*}
        P(X_{t+1}=x' \mid X_t=x, q_t) =
        \frac{\omega_\theta(x',x)}{\sum_{y\in\mathcal N(x)}\omega_\theta(y,x)}
        \ \text{ for } x'\in\mathcal N(x), \quad 0 \text{ otherwise,}
    \end{equation*}
    \begin{equation*}
        \omega_\theta(x',x) =
        \begin{cases}
            \theta & \text{if } d_t(x') \ge d_t(x)
                    \ \text{(the ghost does not get closer to Pacman)},\\
            1 & \text{otherwise.}
        \end{cases}
    \end{equation*}
    We recover \texttt{confused} with $\theta=1$ (uniform over the legal
    moves), \texttt{afraid} with $\theta=2$ and \texttt{scared} with
    $\theta=8$. If a ghost has $a$ moves that do not get it closer to Pacman
    and $b$ moves that do, it flees with probability $\theta a/(\theta a+b)$,
    so a larger $\theta$ means a more fearful ghost. If all the moves belong
    to the same class, the model reduces to a uniform choice.
\end{enumerate}

\section{Implementation}

\begin{enumerate}[label=\alph*.,leftmargin=*]
    \item \textbf{\textit{Leave empty.}}
\end{enumerate}

\section{Experiment}

\begin{enumerate}[label=\alph*.,leftmargin=*]
    \item \textbf{Uncertainty measure.}
    For a ghost with belief $b_t$ (a probability distribution over the free
    cells), we use the Shannon entropy
    \begin{equation*}
        H(b_t) = -\sum_{x} b_t(x)\log_2 b_t(x),
    \end{equation*}
    with the convention $0\log_2 0=0$. It is $0$ when Pacman is certain of the
    ghost position and maximal, $\log_2 |\{\text{free cells}\}|$, when the
    belief is uniform. With several ghosts, we report the value for each ghost
    that has not been eaten yet.

    \item \textbf{Quality measure.}
    Let $x^\star_t$ be the true position of a ghost. We use the expected
    Manhattan distance between the belief and the ground truth,
    \begin{equation*}
        D(b_t, x^\star_t) = \sum_{x} b_t(x)\,
            \big(|x_1-x^\star_{t,1}|+|x_2-x^\star_{t,2}|\big),
    \end{equation*}
    expressed in cells. The lower, the better. Unlike the entropy, it
    penalizes a belief that is concentrated on a wrong location, so the two
    measures are complementary: a confident but wrong filter has a low entropy
    and a large $D$.

    \item \textbf{Results.}
    \textcolor{red}{[TODO: fill in with your own settings and results.]}
    For each layout (\texttt{large\_filter} and \texttt{large\_filter\_walls})
    and each ghost type (\texttt{scared}, \texttt{afraid},
    \texttt{confused}), we ran \textcolor{red}{[N]} independent games with
    different seeds, each lasting \textcolor{red}{[T]} time steps, with one
    ghost and Pacman moving randomly among its legal actions. At every time
    step, the entropy and the distance of the previous questions are recorded
    and then averaged over the trials. The error bars show
    \textcolor{red}{[the standard deviation / the standard error]} over the
    trials. We chose $N$ and $T$ large enough for the curves to be flat and
    the error bars to be stable at the end of the games.

    \begin{figure}[h]
        \centering
        % \includegraphics[width=\linewidth]{results_large_filter.pdf}
        \fbox{\parbox[c][4.2cm][c]{0.9\linewidth}{\centering
            \textcolor{red}{[TODO: plot of entropy and expected distance vs.\
            time step on \texttt{large\_filter}, three ghosts, with error bars]}}}
        \caption{Results on \texttt{large\_filter}.}
    \end{figure}

    \begin{figure}[h]
        \centering
        % \includegraphics[width=\linewidth]{results_large_filter_walls.pdf}
        \fbox{\parbox[c][4.2cm][c]{0.9\linewidth}{\centering
            \textcolor{red}{[TODO: same plots on
            \texttt{large\_filter\_walls}]}}}
        \caption{Results on \texttt{large\_filter\_walls}.}
    \end{figure}

    \item \textbf{Effect of the transition parameter $\theta$.}
    The parameter $\theta$ controls how predictable a ghost is. For
    $\theta=1$ the ghost performs a uniform random walk: the prediction step
    of the filter spreads the belief in all directions and adds a lot of
    uncertainty at each step. When $\theta$ grows, the ghost flees with a
    probability close to $1$ whenever it can, so its motion is nearly
    deterministic and the prediction step adds less uncertainty. We therefore
    expect a lower entropy and a lower distance for \texttt{scared} than for
    \texttt{afraid}, and for \texttt{afraid} than for \texttt{confused}. The
    layout also matters: in \texttt{large\_filter\_walls} the walls reduce the
    number of reachable cells and the branching factor of the ghost, which
    both lower the uncertainty added by the prediction step, whereas in the open
    layout the belief can spread in more directions. A fleeing ghost may also
    get trapped in a dead end or a corner, where the belief concentrates.
    \textcolor{red}{[TODO: confirm or correct these statements with
    the values observed in the figures, with default sensor variance.]}

    \item \textbf{Effect of the sensor variance.}
    The sensor variance $\sigma^2$ sets $n=4\sigma^2$, and hence the width
    of the sensor likelihood: a measurement only tells that the ghost is on a
    ring of thickness proportional to $\sigma$ around Pacman. The larger
    $\sigma^2$, the thicker the ring, the flatter the posterior after the
    update step, and the larger the entropy and the distance of the belief;
    convergence also takes longer since more measurements are needed to
    discard cells. For very small variances ($\sigma^2<0.25$ gives $n=0$) the
    distance is exact and the posterior lies on a thin ring, but the ghost
    position along this ring remains ambiguous until Pacman has moved
    and the rings of several time steps can be combined.
    \textcolor{red}{[TODO: back it up with the values observed for
    several values of \texttt{--sensorvariance}.]}

    \item \textbf{Controller to eat ghosts.}
    Pacman only knows its position, its legal actions and the belief of each
    ghost. A simple controller is: (i) pick a target among the ghosts that have
    not been eaten, for instance the one whose belief mass is the most
    concentrated or the closest to Pacman; (ii) take as goal the most likely cell
    $\arg\max_x b(x)$ of this belief; (iii) among the legal actions,
    choose the one leading to the cell closest to the goal. To be more robust, the
    goal can be replaced by the cell minimizing the expected distance
    $\sum_x b(x)\,d(\cdot,x)$ to the belief, and the distance by the
    shortest path distance in the maze (computed with a breadth-first search on
    the grid of free cells), so that Pacman does not get stuck against a wall.
    When the belief is still diffuse, the controller can also move toward the
    region that reduces the entropy the most before chasing.

    \item \textbf{\textit{Leave empty.}}
\end{enumerate}

% ==============================================================================

\end{document}